\chapter{Developer Contribution Guide}
\label{ch:contributing}

This chapter outlines the guidelines for developers contributing to the RGUKT Connect codebase, covering branching strategies, coding standards, and deployment tests.

\section{Branching Strategy \& Workflow}
The repository uses a Git Flow branching model:
\begin{itemize}
    \item \textbf{`main`:} Production branch. Contains stable, tested releases. Direct commits to `main` are restricted.
    \item \textbf{`develop`:} Integration branch. All features and bug fixes are integrated here first.
    \item \textbf{`feature/*`:} Feature development branches. Created from `develop` and merged back via Pull Requests (PRs).
    \item \textbf{`hotfix/*`:} Urgent production fixes. Created from `main` and merged into both `main` and `develop`.
\end{itemize}

\section{Coding Standards \& Formatting}
Developers must follow consistent formatting standards across layers.

\subsection{Backend Java Code Standards}
\begin{itemize}
    \item Use standard camelCase naming for variables and methods, and PascalCase for classes.
    \item Maintain a clean separation of concerns: Controller classes must not execute business logic directly, and service classes must not interact with database JDBC layers directly.
    \item Document public API endpoints using Swagger or SpringDoc annotations.
\end{itemize}

\subsection{Frontend React Code Standards}
\begin{itemize}
    \item Write functional components using standard ES6 syntax and React Hooks.
    \item Manage local states with the `useState` hook, and use custom Context hooks for shared global state.
    \item Enforce styling rules using Tailwind CSS utility classes, and extract complex layout blocks into reusable UI components.
\end{itemize}

\section{Commit Message Conventions}
Commits should use structured, semantic prefixes:
\begin{itemize}
    \item \textbf{`feat:`} A new functional capability (e.g., `feat: integrate S3 avatar upload`).
    \item \textbf{`fix:`} A bug fix (e.g., `fix: resolve CORS header mismatch`).
    \item \textbf{`docs:`} Documentation updates (e.g., `docs: add deployment instructions`).
    \item \textbf{`refactor:`} Code refactoring that does not change functionality (e.g., `refactor: optimize JPA fetch join`).
    \item \textbf{`test:`} Adding or modifying tests (e.g., `test: add user service unit tests`).
\end{itemize}

\section{Pre-PR Checklist}
Before opening a Pull Request, developers must complete the following steps:
\begin{enumerate}
    \item Ensure the application compiles without errors:
    \begin{lstlisting}[language=bash]
    # Backend check
    mvn clean compile
    
    # Frontend check
    npm run build
    \end{lstlisting}
    \item Run all context and unit tests:
    \begin{lstlisting}[language=bash]
    mvn test
    \end{lstlisting}
    \item Verify formatting and linting:
    \begin{lstlisting}[language=bash]
    npm run lint
    \end{lstlisting}
    \item Test API endpoints using local REST clients (such as Postman or Bruno).
\end{enumerate}
